\Needspace{13\baselineskip}
\chapterrule
\section{Attention：从直接读取到 Q/K/V 分工}
\label{intro:21c-naive-attention}
\label{intro:21b-why-attention}
\phantomsection\label{concept:attention}

\begin{idea}
\noindent\emph{按相关性读取上下文}\par

Attention（注意力）是一种按相关性读取信息的方法：
当前位置先对\textbf{可读取的每个位置}算一个匹配分数，
再把这些分数变成\textbf{权重}，
最后把各位置提供的信息按权重\textbf{加起来}。

它最核心的三个动作就是
\textbf{比较 $\to$ 分配权重 $\to$ 加权读取}。
\end{idea}

\subsection{从逐步传递到直接读取}

\conceptlink{rnn}{RNN} 让模型按顺序读，前文要经过一份份状态才能传到当前位置。
先比较两条信息路径：
% Side-by-side panels keep labels readable while fitting below the opening text.
\par\noindent\begin{minipage}{\linewidth}
\tagstructbegin{tag=Figure,alttext={左图 RNN 依次从春天到花再到开了传递状态；右图 Attention 从三个可见位置直接读取到位置三。右图线宽示意不同读取权重。}}
\begin{tcolorbox}[enhanced,colback=paper,colframe=linegray,
 boxrule=.7pt,arc=2mm,left=3mm,right=3mm,top=2mm,bottom=2mm]
\centering
\begin{tikzpicture}[>=Latex,x=1mm,y=1mm,
 word/.style={draw=linegray,fill=white,rounded corners=2pt,
  minimum width=17mm,minimum height=9mm,align=center,font=\footnotesize},
 output/.style={draw=flowcolor,fill=paper,rounded corners=2pt,
  minimum width=22mm,minimum height=10mm,align=center,font=\footnotesize}]
\node[font=\sffamily\bfseries\small,anchor=west] at (0,25) {RNN：依次传递};
\node[word] (r1) at (9,12) {春天\\$h_1$};
\node[word] (r2) at (30,12) {花\\$h_2$};
\node[word] (r3) at (51,12) {开了\\$h_3$};
\draw[->,very thick,flowcolor] (r1)--(r2);
\draw[->,very thick,flowcolor] (r2)--(r3);
\node[font=\footnotesize,align=center] at (30,-2) {较早信息经过中间状态};
\draw[linegray] (68,28)--(68,-9);
\node[font=\sffamily\bfseries\small,anchor=west] at (75,25) {Attention：直接读取};
\node[word] (a1) at (84,13) {春天\\位置 1 输入};
\node[word] (a2) at (108,13) {花\\位置 2 输入};
\node[word] (a3) at (132,13) {开了\\位置 3 输入};
\node[output] (at) at (108,-4) {位置 3 的\\读取结果};
\draw[->,flowcolor,line width=2pt] (a1.south) to[out=-80,in=165] (at.west);
\draw[->,flowcolor,line width=.6pt] (a2.south)--(at.north);
\draw[->,flowcolor,line width=1.15pt] (a3.south) to[out=-100,in=15] (at.east);
\end{tikzpicture}
\par
\end{tcolorbox}
\begin{center}\captionof{figure}{RNN 逐步传递与 Attention 直接读取的对照；线宽示意读取权重}\label{fig:attention-rnn}\end{center}
\tagstructend
\end{minipage}\par

RNN 的第 3 步能读到\textbf{上一份状态} $h_2$；
想重点利用第 1 个词的信息，就要靠这份状态此前保留了相关内容。

换一种设计：读到第三个词“开了”时，
让第三个位置直接查看位置 1、2、3，各自决定读多少。
图中右侧的线宽表示读取比例不同。接下来要解决的，就是怎样算出这些比例。

\subsection{从平均读取到按需读取}

设当前位置可以读取三个位置，它们各自提供一条\textbf{内容向量}（行向量）：
\[
\vect v_1^\top=(1,0),\qquad
\vect v_2^\top=(0,1),\qquad
\vect v_3^\top=(1,1).
\]
最简单的方法是把它们\textbf{平均}：
\[
\vect y_{\rm avg}^\top
=\tfrac13\bigl(\vect v_1^\top+\vect v_2^\top+\vect v_3^\top\bigr)
=(2/3,\ 2/3).
\]
无论当前位置想查找什么，只要内容不变，平均结果就不变。

改为给三条内容分配权重 $a_1,a_2,a_3$，
并规定权重非负、总和为 $1$：
\begin{equation}
\boxed{\vect y^\top=a_1\vect v_1^\top+a_2\vect v_2^\top+a_3\vect v_3^\top,\qquad
a_i\ge0,\quad \sum_{i=1}^3 a_i=1.}
\label{eq:attention-design-goal}
\end{equation}
平均对应 $a_1=a_2=a_3=1/3$。
要让读取结果随当前位置的需要而变化，权重就得取决于
\textbf{当前想找什么}与\textbf{各位置凭什么被找到}的匹配程度。
于是每个位置在读取中承担三种职责：发起比较、被比较、交出内容。
本章先让同一条输入向量兼任这三种职责。

\subsection{距离变长，同一条线索还剩多少？}
\label{experiment:distance-reading}

下面做一个只比较“距离”的条件实验。先把要保留的信息缩成一个数：
第一个位置的输入为 $1$，之后的位置都为 $0$；当前位置在第 $d+1$ 个位置，
距第一位置 $d$ 步。以第一个位置为读取目标，手设匹配分数
$s_1=\ln9$，其余位置的分数为 $s_2=\cdots=s_{d+1}=0$。
\textbf{固定上述匹配分数，不训练 $Q、K$，只改变距离，比较两种路径保留的信号。}

\textbf{沿状态传递：}沿用上一章的 $h_t=x_t+0.5h_{t-1}$、$h_0=0$，
第一个位置得到 $h_1=1$；之后每步都乘 $0.5$。
例如距离为 $2$ 时，$h_3=0.5\times0.5\times1=0.25$；一般为 $h_{d+1}=0.5^d$。

\textbf{直接加权读取：}保留每个位置的内容，令 $v_1=1$，其余 $d$ 个内容均为 $0$。
在上述手设分数下，用已学的 Softmax 分配权重，第一个位置的权重和读出为
\[
a_1=\frac{e^{\ln9}}{e^{\ln9}+d\,e^0}=\frac{9}{9+d},\qquad
y=a_1\times1+\sum_{i=2}^{d+1}a_i\times0=\frac{9}{9+d}.
\]
\begin{center}\small
\renewcommand{\arraystretch}{1.2}
\tagpdfsetup{table/header-rows={1}}
\begin{tabularx}{\linewidth}{@{}p{24mm}XX@{}}
\toprule
距离 $d$ & RNN：状态中的信号 $0.5^d$ & 直接读取（手设分数）：$9/(9+d)$\\\midrule
$1$ & $0.5000$ & $0.9000$\\
$2$ & $0.2500$ & $0.8182$\\
$4$ & $0.0625$ & $0.6923$\\
$8$ & $0.0039$ & $0.5294$\\
\bottomrule
\end{tabularx}
\end{center}
相隔 $8$ 步时，这个 RNN 只剩约 $0.004$ 的信号，直接读取仍约为 $0.529$。
前者反复缩小同一份状态；后者在原位置读取内容，却也会因候选增多而\textbf{稀释权重}。

这张表比较的是两条手设规则下的\textbf{衰减与稀释}。
把递推系数改成 $1$，并保持后续输入全为零，就能一直保留这个 $1$。
要理解长历史的困难，还需检查状态装得下什么，以及训练信号经过多少次更新。

\subsection{定长状态：保存一个数与保存多条绑定不同}

\textbf{容量与覆盖：}RNN 把长度不断增长的历史压进固定 $d_h$ 维状态。
在固定数值精度下，这份状态能区分的内容有限；新输入与旧信息共用状态，
更新时可能混合或覆盖此前的记录。
例如将系数改成 $1$ 后，规则 $h_t=x_t+h_{t-1}$
会把 $(1,0)$ 与 $(0,1)$ 都压成最终状态 $1$。
若最后才问“第一项是多少”，两个答案应为 $1$ 和 $0$，仅凭这个和已无法还原。
更宽的状态或 LSTM、GRU 的门控可以学会保留更多信息，但仍要选择保存哪些内容。

\textbf{梯度路径：}这个线性递推对早期状态的导数也连乘 $d$ 次系数，
$\partial h_{d+1}/\partial h_1=a^d$。
一般 RNN 则连乘各步状态更新的导数矩阵，可能使梯度缩小或放大；
依赖距离越长，递推路径越长，训练就越需要处理这种数值问题。
直接注意力在一层内连接查询位置与被读取位置，提供较短的跨位置路径；
它付出的代价是保存多条记录，并计算匹配。
递推训练困难与路径比较分别见\hyperref[source:rnn-training]{RNN 训练论文}
和\hyperref[source:attention-original]{Transformer 原论文}。

\Needspace{12\baselineskip}
\subsection{当场检验：两条记录、两个查询}
\label{experiment:record-query-handcalc}

设记录甲为“花：$+1$，窗：$-1$”，记录乙把两个值交换。
分别查询“花”和“窗”。这里直接编码记录名：花为 $(1,0)$、窗为 $(0,1)$，
不经过主线词表或 Embedding；两份记录的名字与排列相同，只有值不同。
把两个查询和两个记录名都按“花、窗”排成 $U=I_2$，
固定匹配规则为名字向量点积乘 $\ln3$，不训练参数：
\[
S=(\ln3)UU^\top=\begin{bmatrix}\ln3&0\\0&\ln3\end{bmatrix},\qquad
A=\softmax_{\rm row}(S)=\begin{bmatrix}3/4&1/4\\1/4&3/4\end{bmatrix}.
\]
$A:2\times2$ 的行对应查询，列对应记录；同名记录拿到 $3/4$ 的权重。
两份记录各自的值都是 $2\times1$ 的列，读出为
\[
V_{\text{甲}}=\begin{bmatrix}1\\-1\end{bmatrix},\quad
AV_{\text{甲}}=\begin{bmatrix}1/2\\-1/2\end{bmatrix};\qquad
V_{\text{乙}}=\begin{bmatrix}-1\\1\end{bmatrix},\quad
AV_{\text{乙}}=\begin{bmatrix}-1/2\\1/2\end{bmatrix}.
\]
例如查询“花”时，甲读出 $(3/4)\times1+(1/4)\times(-1)=1/2$，乙读出 $-1/2$；
按正负判别，两份记录的答案都能区分。
均匀读取的矩阵每项都是 $1/2$，乘任一份值都得到 $(0,0)^\top$。
选择性读取保住了答案的符号；有限 Softmax 权重仍混入另一条记录，读出幅度是原值的一半。

接收完整历史的 RNN 或固定窗口 MLP 也能解决这个小任务。
这里的匹配来自固定名字编码；怎样通过训练学到匹配，见第~\ref*{intro:22c-language-training}~章
\hyperref[experiment:attention-selection]{记录查询训练对照}（第~\pageref*{experiment:attention-selection}~页）的均匀读取与学习 Q/K 结果。
下面用三个二维输入继续拆解匹配分数的来源。

\subsection{分数从哪里来？三个二维输入的零参数读取}

前面的距离实验直接指定了匹配分数。
现在回到前面用于平均读取的三条二维输入，展示一种最简单的分数来源：
$\vect x_1^\top=(1,0)$，$\vect x_2^\top=(0,1)$，$\vect x_3^\top=(1,1)$。
每条输入既拿去比较，也被比较，还作为读出的内容，
也就是暂时令 $Q=K=V=X$，三种职责共用输入，不设置可训练投影。
后面会沿用同一组输入，用三张可训练投影矩阵分别承担三种职责。
先将所有点积写成矩阵：
\[
X=\begin{bmatrix}1&0\\0&1\\1&1\end{bmatrix}
\ (3\times2),\quad
\underbrace{S}_{3\times3}=XX^\top
=\boxed{\begin{bmatrix}1&0&1\\0&1&1\\1&1&2\end{bmatrix}}.
\]
此处\textbf{未缩放、未加因果掩码}。第 $i$ 行是位置 $i$ 与三个位置的点积。
例如第三行 $(1,1,2)$：只对这条 Query 的三个分数计算 Softmax，得
$(0.212,0.212,0.576)$，继续把输入作为内容向量，
读出约为 $(0.788,0.788)$。
与平均读取的 $(0.667,0.667)$ 相比，结果偏向了匹配分数最高的位置 3。
不增加任何参数，也能完成一次“比较、分配、读取”。

\Needspace{14\baselineskip}
\begin{boundarybox}
\textbf{零参数点积的匹配规则是固定的：}
(1) \textbf{分数对称。}$S=XX^\top=S^\top$：
“位置 $i$ 拿自己去比 $j$”与“位置 $j$ 拿自己去比 $i$”得到同一个原始分数。
(2) \textbf{长度也计入分数。}点积同时受方向和长度影响。
$\vect x_3^\top=(1,1)$ 最长，在三行里都拿到最高或并列最高的分数：
第一行 $(1,0,1)$、第二行 $(0,1,1)$、第三行 $(1,1,2)$。
逐行做 Softmax 后，第一行对位置 $1,3$ 给出并列最大权重；
第二行对位置 $2,3$ 给出并列最大权重；只有第三行对位置 $3$ 给出唯一最大权重。
这些分配由当前输入的方向、长度与固定的点积规则共同决定，尚未通过训练按任务需要调整。
两点同出一源：这一版本没有可调的数，训练无法让模型学会“按哪些特征匹配”。
\end{boundarybox}

要放宽这些限制，就让“拿什么询问”“凭什么被找到”“读出什么内容”
各由一组可训练投影产生。
下面就把三种职责拆开，并沿用同一组输入比较读出变化。



\begin{idea}
\noindent\emph{同一个位置，三种身份}\par

前面的零参数版本让每条输入向量同时负责发起比较、被比较和交出内容，
匹配规则因此固定：原始分数对称，匹配同时受输入方向和长度影响。
现在给这三种职责各配一组\textbf{可训练投影}：
同一条输入分别变成 Query（查询）、Key（键）和 Value（值）。
比较、分配权重、加权读取这三个动作不变，变的是参与比较和读取的向量。
\end{idea}

\subsection{Query、Key、Value 各负责什么？}
\phantomsection\label{concept:qkv}

在点积 Attention 中，三者分工明确：
\begin{itemize}
\item \textbf{Query（查询，记作 $\vect q^\top$）：}
当前这个位置\textbf{拿什么条件来查}？它用于发起匹配。
\item \textbf{Key（键，记作 $\vect k_i^\top$）：}
第 $i$ 个位置\textbf{凭什么被找到}？它与 Query 比较，
产生一个\textbf{匹配分数}。
\item \textbf{Value（值，记作 $\vect v_i^\top$）：}
第 $i$ 个位置\textbf{被找到后交出什么内容}？
它不是匹配分数，而是最终\textbf{被加权相加}的向量。
\end{itemize}

\begin{center}\QKVlegend\end{center}
Query 和 Key 用来计算权重，Value 则是按权重组合的内容。

% The Figure structure intentionally includes the non-floating caption.
\tagstructbegin{tag=Figure,alttext={Query 匹配 Key，按得到的权重对 Value 加权求和}}
\begin{tcolorbox}[enhanced,colback=paper,colframe=linegray,
  boxrule=.7pt,arc=2mm,left=3mm,right=3mm,top=3mm,bottom=3mm,
  fonttitle=\sffamily\bfseries,coltitle=ink,
  title={从“按需要读取”到 Q/K/V：每个字母都有工作}]
\small
\begin{center}
\begin{tikzpicture}[>=Latex,
  b/.style={draw=linegray,fill=white,rounded corners=3pt,
             minimum width=52mm,minimum height=12mm,align=center,font=\small},
  head/.style={draw=blue,fill=blueback,rounded corners=3pt,
             minimum width=118mm,minimum height=10mm,align=center,font=\small},
  end/.style={draw=teal,fill=greenback,rounded corners=3pt,
             minimum width=118mm,minimum height=11mm,align=center,font=\small}]
\node[head] (goal) at (0,1.95) {目标：让当前位置按相关性读取三个位置的信息};
\node[b] (find) at (-3.15,.6) {先决定读谁、读多少};
\node[b] (read) at (3.15,.6) {再决定实际取回什么};
\node[b] (match) at (-3.15,-.85)
  {\textcolor{blue}{Query} $\vect q_3$ 与 \textcolor{teal}{Key} $\vect k_i$\\点积得到分数 $s_i$};
\node[b] (value) at (3.15,-.85)
  {每个位置提供 \textcolor{orange}{Value} $\vect v_i$\\它是被读取的内容};
\node[b] (weight) at (-3.15,-2.35) {Softmax：分数 $\to$ 权重 $a_i$};
\node[end] (out) at (0,-3.85)
  {结果：$\displaystyle\vect y_3=\sum_{i=1}^{3}a_i\vect v_i$};
\draw[->,blue,thick] (goal.south) to[out=-115,in=90] (find.north);
\draw[->,teal,thick] (goal.south) to[out=-65,in=90] (read.north);
\draw[->,blue,thick] (find)--(match);
\draw[->,orange,thick] (read)--(value);
\draw[->,blue,thick] (match)--(weight);
\draw[->,blue,thick] (weight.south) to[out=-90,in=150] (out.north west);
\draw[->,orange,thick] (value.south) to[out=-90,in=30] (out.north east);
\end{tikzpicture}
\end{center}
\textbf{注意：}这里 $i=1,2,3$ 是\textbf{候选位置编号}。
第三个位置既能发出 Query，也会提供自己的 Key 和 Value。
Q/K/V 是\textbf{每个位置参与计算时扮演的角色}，
不是把三个词分别命名为 Q、K、V。
\end{tcolorbox}
\begin{center}\captionof{figure}{Query、Key 与 Value 的分工：先决定匹配权重，再读取内容}\label{fig:why-qkv}\end{center}
\tagstructend


\textbf{Q、K、V 是三种角色。}在 Self-Attention 中，
第 3 个 Token 一方面能发出自己的 Query，
另一方面也像其他位置一样提供 Key 和 Value。
每个位置原本有一条输入表示 $\vect x_i^\top$，
在真实网络中可以经过\textbf{不同的可训练投影}，得到
\[
\vect q_i^\top=\vect x_i^\top W_Q,\qquad
\vect k_i^\top=\vect x_i^\top W_K,\qquad
\vect v_i^\top=\vect x_i^\top W_V.
\]
这里约定每个 Token 的输入 $\vect x_i^\top$ 是\textbf{行向量}，
$W_Q:d_{\rm model}\times d_k$；行列写法的对应关系见\sectionentry{sec:layout-matrix-card}。
三组投影参数都可以通过训练学习。

为什么区分这三种向量？因为\textbf{用于判断相关性的特征}与\textbf{希望读出的内容}可以不同，
Value 的维度也不必等于 Key 的维度。
独立投影让网络能分别学习匹配和读出的表示；简化模型也可以让部分投影相同。

打个类比：查资料时，问题像 Query，索引像 Key，正文内容像 Value。

先看\textbf{一次查询}：
位置 3 发出 Query，分别与位置 1、2、3 的 Key 匹配，
得到三条权重，再从这三个位置的 Value 读取内容。

\subsection{匹配：独立投影让两个方向的分数可以不同}

沿用前面三条二维输入组成的 $X$，每行一个位置；
为本次前向计算指定三张投影矩阵的数值：
\[
X=\begin{bmatrix}1&0\\0&1\\1&1\end{bmatrix},\qquad
W_Q=I_2,\qquad
W_K=\begin{bmatrix}1&0\\1&1\end{bmatrix},\qquad W_V=I_2.
\]
先直接用点积作为分数，并让每个位置都可读取三个位置。矩阵乘法得到
\[
Q=XW_Q=X,\quad
K=XW_K=\begin{bmatrix}1&0\\1&1\\2&1\end{bmatrix},\quad
V=XW_V=X.
\]
Key 的三行变成 $(1,0)$、$(1,1)$、$(2,1)$，不再等于输入（只有第一行碰巧相同）。
三条 Query 与三条 Key 两两做点积，得到 $3\times3$ 的分数矩阵：
\begin{equation}
S=QK^\top
=\boxed{\begin{bmatrix}1&1&2\\0&1&1\\1&2&3\end{bmatrix}}.
\label{eq:asymmetric-qk-example}
\end{equation}
\textbf{对照零参数版本：}原始 $XX^\top$ 必定对称，
而此处 $S_{12}=\boxed{1}$、$S_{21}=\boxed{0}$。
位置 1 查询位置 2 与位置 2 查询位置 1 的匹配分数可以不同；
这正是\textbf{拆开 Query 和 Key 的新增表达能力}。
判断匹配是否对称要看原始分数 $S$：逐行 Softmax 各用自己的分母，
权重矩阵是否对称反映不了 $S$ 是否对称。
把 $S$ 逐行做 Softmax 后得到的\textbf{权重矩阵记作 $A$}；
$A_{ij}$ 表示位置 $i$ 读取位置 $j$ 的权重。
例如回到零参数版本，$S_{13}=S_{31}=1$，但
$A_{13}=e/(2e+1)\approx0.4223$，$A_{31}=e/(2e+e^2)\approx0.2119$。
两行分母不同，同样的原始分数也会得到不同权重。

\subsection{分配：以第三个位置为 Query 计算权重}

当前第三行 Query 是 $\vect q^\top=(1,1)$，
与三个 Key 的分数是 $(1,2,3)$：
\[
\vect a^\top=\softmax(1,2,3)
=\frac{(e,e^2,e^3)}{e+e^2+e^3}
\approx\boxed{(0.090,0.245,0.665)}.
\]
预测下一 Token 时，当前位置只读自己和前面的 Token。
第三个位置的三个 Key 均可见；第一、二位置还要排除各自的未来位置，
\hyperref[concept:causal-mask]{下一章}会算出这两行的因果版本。
本节使用\textbf{未缩放}的原始点积以便比较投影；
缩放的目的与算法在\sectionentry{concept:scaled-dot-product}中再解释。

\subsection{读出：Value 与 Key 分开设计}

由于 $W_V=I_2$，三行 Value 仍是
$(1,0)$、$(0,1)$、$(1,1)$。
\begin{example}
\[
\vect y^\top=\vect a^\top V
\approx(0.090,0.245,0.665)
\begin{bmatrix}1&0\\0&1\\1&1\end{bmatrix}
=\boxed{(0.755,0.910)}.
\]
这与前面未投影时的 $(0.788,0.788)$ 不同：
\textbf{输入和 Value 没变，只改变 Key 的投影就改变了读取内容的比例。}
\end{example}

如果只把第一行 Value 改成 $(0,2)$，
分数和权重完全不变，输出则约为
$(0.665,\;0.090\times2+0.245+0.665)=(0.665,1.090)$。
所以 Key 决定匹配、Value 决定取出的内容，
两者的作用可以分别核对。

% The Figure structure intentionally includes the non-floating caption.
\tagstructbegin{tag=Figure,alttext={二维独立例中位置三的 Query 为一一，位置一到三的 Key 匹配分数依次为一、二、三；Softmax 权重为零点零九零、零点二四五、零点六六五，Value 加权结果约为零点七五五、零点九一零。}}
\begin{tcolorbox}[enhanced,colback=paper,colframe=linegray,
boxrule=.7pt,arc=2mm,left=3mm,right=3mm,top=3mm,bottom=3mm,
coltitle=ink,fonttitle=\sffamily\bfseries,
title={三条 Key 的匹配分数与读取权重}]
\small
独立例：位置 3 的 Query 为 $(1,1)$；Key 经独立投影，Value 仍是输入。
\begin{center}
\begin{tikzpicture}[x=1cm,y=.9cm,>=Latex]
\node[anchor=west,font=\small\bfseries] at (0,.5) {位置};
\node[anchor=west,font=\small\bfseries] at (2.6,.5) {点积分数};
\node[anchor=west,font=\small\bfseries] at (5.1,.5) {Softmax 权重};
\node[anchor=west,font=\small\bfseries] at (10.5,.5) {Value};
\foreach \yy/\word/\s/\weight/\len/\vv in {
 0/位置 1/1/0.090/.38/{(1,0)},
 -1.0/位置 2/2/0.245/1.03/{(0,1)},
 -2.0/位置 3/3/0.665/2.79/{(1,1)}
}{
 \node[anchor=west,font=\small] at (.1,\yy) {\word};
 \node[font=\small] at (3.5,\yy) {$\s$};
 \fill[paper] (5.1,\yy-.18) rectangle (8.1,\yy+.18);
 \fill[datacolor] (5.1,\yy-.18) rectangle ({5.1+\len},\yy+.18);
 \draw[linegray] (5.1,\yy-.18) rectangle (8.1,\yy+.18);
 \node[anchor=west,font=\small] at (8.25,\yy) {$\weight$};
 \node[anchor=west,font=\small] at (10.65,\yy) {$\vv$};
}
\node[anchor=west,font=\small] at (.1,-2.9)
 {$0.090(1,0)+0.245(0,1)+0.665(1,1)\approx(0.755,0.910)$};
\end{tikzpicture}
\end{center}
\end{tcolorbox}

\begin{center}\captionof{figure}{独立 Key 投影下的第三行 Query：匹配、归一化与 Value 读出}\label{fig:attention-weights}\end{center}
\tagstructend


\subsection{Attention 与 RNN 怎样传递信息？}

\begin{center}
\small
\renewcommand{\arraystretch}{1.35}
\tagpdfsetup{table/header-rows={1}}
\begin{tabularx}{\linewidth}{@{}p{33mm}XX@{}}
\toprule
比较的问题 & RNN & 本章的 Attention\\
\midrule
当前位置怎样利用前文？ & 接收上一时刻传来的状态 & 直接对多个候选位置匹配和读取\\
输出是什么？ & 更新后的状态向量 & 多条 Value 的加权组合向量\\
是否需要逐步传状态？ & 是 & 核心读取运算不需要这样递推\\
是否使用可训练参数？ & 可以，各时间步共享更新参数 & Q/K/V 的投影在真实模型中可学习\\
\bottomrule
\end{tabularx}
\end{center}

回到主线，还需要把两条规则放进模型：\textbf{因果掩码}
让当前位置只读取已经可见的 Token；序列变长时，
注意力的匹配与缓存也要支付相应成本。
本章先算一个 Query，下章把所有行的 Softmax 和 Value 读取合成矩阵运算，并加上因果掩码。

前面在固定投影下算出了读取权重。接下来还要让投影通过训练适应任务：
\textbf{能读到前文，不等于会挑选前文。}
\hyperref[experiment:attention-selection]{第~\ref*{experiment:attention-selection}~节的记录查询对照}
将检验 Q/K 怎样学会按查询选择记录。

\chapterchecktasks{21c-naive-attention}{理解、手算与判断}{
\label{check:21b-why-attention}
\item\textbf{手算：}本章 $\vect x_1^\top=(1,0),\vect x_2^\top=(0,1),\vect x_3^\top=(1,1)$。计算 $S=XX^\top$ 的第一行和第三行，标出 $S_{13}$ 与 $S_{31}$；再按三行分数判断各行哪些位置得到最大权重，注意并列。\chapterinlinehint{让对应输入行分别与所有输入行做点积，再比较同一行的分数。}
\item\textbf{预测：}对每一行分开做 Softmax 后，$A_{13}$ 与 $A_{31}$ 还一定相等吗？请先比较两行的分母。\chapterinlinehint{将两个相同分子与各自行的分母对照。}
\item\textbf{找错：}同学说“直接点积能自定义‘我如何问’和‘别人如何被问’这两套不同视角”。这一零参数方案缺少什么？\chapterinlinehint{检查提问与匹配是否使用两套可训练矩阵。}
\item\textbf{定位：}一句话说明 Query、Key、Value 三者在“计算匹配／提供内容”中的分工，并标出分数和读出的形状：$\vect q^\top:1\times2,K:3\times2,V:3\times2$。\chapterinlinehint{分别追踪分数路径与内容加权路径。}
\item\textbf{手算：}只有两个 Value 为 $(1,0),(0,1)$，权重为 $(0.2,0.8)$，求加权读出。\chapterinlinehint{把每个 Value 乘其权重，再逐分量相加。}
\item\textbf{找错：}同学说“保持 Q、K 不变，修改第三条 Value 后，每条原始匹配分数必定改变”。指出混淆了哪两个路径。\chapterinlinehint{写出分数公式，检查其中是否出现 Value。}
}


\par\smallskip
同一组输入在拆分职责后读出不同结果；可训练投影使匹配方式也能由任务反馈调整。
$XX^\top$ 必定对称来自 $Q=K=X$ 的特殊设定；一般的 $QK^\top$ 不要求对称。
下一章把读取组织成矩阵，一次算完所有位置，并加入因果掩码。
